\section{Validation et résultats}
\label{sec:resultats}

Les résultats portent sur l'OAT 4,10~\% 2046, pour des options d'échéance 1, 2, 5, 7 et 10 ans et des strikes
clean de 90 à 110~\% du forward clean, avec les paramètres retenus aux sections~\ref{sec:taux}
et~\ref{sec:spread}. On valide d'abord le prix exact, puis on reprend les hypothèses H5, H3 et H4.

\subsection{Validation du prix exact et précision des poids gelés}

Le prix~\eqref{eq:prix_exact} est exact au calcul de l'intégrale près. On le valide par un Monte-Carlo
indépendant. Le prix s'écrit
$\E^{\Q}\big[e^{-\int_0^{T_0}r}\,(\bar B(T_0)-X)^+\big]$, et tout ce qui y entre est fonction du vecteur
gaussien $Z = \big(x(T_0), y(T_0), \int_0^{T_0}x, \int_0^{T_0}y\big)$, de moyenne nulle sous $\Q$ et de covariance
donnée en annexe~A. On simule $Z$ exactement, sans discrétisation en temps : le facteur d'actualisation vaut
$P^M(0,T_0)\,e^{-\frac12V_x(0,T_0)}\,e^{-\int_0^{T_0}x}$, et $\bar B(T_0)$ s'obtient par la formule
fermée~\eqref{eq:zc_oat}, le facteur spread étant décalé du recentrage $\delta$. La simulation utilise 400\,000
trajectoires avec variables antithétiques.

Trois contrôles précèdent le pricing, à cinq ans. La moyenne simulée de $e^{-\int r}$ vaut 0,863309, pour
$P^M(0,5) = 0{,}863309$. Celle de $e^{-\int \bar r}$ vaut 0,839610, pour $\bar P^M(0,5) = 0{,}839609$, ce qui
teste le recalage de la courbe OAT et toute la matrice de covariance. Enfin, le forward simulé vaut 87,903, pour
un forward repo de 87,908 : la loi forward et le recentrage sont corrects. Dans le cas limite à un facteur
($\sigma_y = 0$, courbe OAT confondue avec la courbe €STR), le prix exact du call à la monnaie à cinq ans,
6,254832, est égal à celui de Jamshidian à la sixième décimale.

\begin{table}[htbp]
\centering\small
\caption{Prix exact, Monte-Carlo (MC, demi-largeur de l'intervalle de confiance à 95~\%) et approximation à poids
gelés, $T_0 = 5$ ans (prix pour 100 de nominal ; écarts relatifs en \%)}
\label{tab:validation}
\begin{tabular}{@{}llcccccc@{}}
\toprule
Strike & Type & Exact & MC & IC 95~\% & Exact/MC $-1$ & Poids gelés & Gelés/exact $-1$ \\
\midrule
90~\%  & call & 9,9609 & 9,9513 & 0,0249 & 0,10 & 10,0170 & 0,56 \\
90~\%  & put  & 2,4752 & 2,4696 & 0,0118 & 0,23 & 2,5313 & 2,27 \\
100~\% & call & 5,7627 & 5,7512 & 0,0259 & 0,20 & 5,7695 & 0,12 \\
100~\% & put  & 5,7627 & 5,7552 & 0,0131 & 0,13 & 5,7695 & 0,12 \\
110~\% & call & 3,0727 & 3,0661 & 0,0217 & 0,22 & 3,0292 & $-1{,}42$ \\
110~\% & put  & 10,5584 & 10,5558 & 0,0085 & 0,02 & 10,5149 & $-0{,}41$ \\
\bottomrule
\end{tabular}
\end{table}

Le prix exact reste dans l'intervalle de confiance du Monte-Carlo pour tous les strikes
(tableau~\ref{tab:validation}) et, à la monnaie, à toutes les échéances, où la demi-largeur relative de
l'intervalle va de 0,4 à 0,5~\%. Le prix exact sert donc de référence pour mesurer l'erreur de l'approximation à
poids gelés, sans bruit de simulation.

\begin{table}[htbp]
\centering\small
\caption{Erreur relative de l'approximation à poids gelés (poids forward recentrés) contre le prix exact, en \%}
\label{tab:poids_geles}
\begin{tabular}{@{}lccccc@{}}
\toprule
 & 1 an & 2 ans & 5 ans & 7 ans & 10 ans \\
\midrule
Call 90~\%  & 0,29 & 0,46 & 0,56 & 0,52 & 0,38 \\
Call 100~\% & 0,04 & 0,08 & 0,12 & 0,11 & 0,08 \\
Call 110~\% & $-2{,}98$ & $-2{,}18$ & $-1{,}42$ & $-1{,}17$ & $-0{,}92$ \\
Put 90~\%   & 4,28 & 3,22 & 2,27 & 1,91 & 1,45 \\
Put 110~\%  & $-0{,}28$ & $-0{,}39$ & $-0{,}41$ & $-0{,}37$ & $-0{,}28$ \\
\midrule
\multicolumn{6}{@{}l}{\footnotesize Poids gelés à leur valeur d'aujourd'hui, call 100~\% : de $-0{,}04$ à $+0{,}50$ (à 10 ans)} \\
\bottomrule
\end{tabular}
\end{table}

À la monnaie, l'approximation de \citet{Russo} s'écarte du prix exact de 0,12~\% au plus lorsque les poids sont
gelés à leur valeur forward recentrée, contre 0,5~\% à dix ans s'ils sont gelés à leur valeur d'aujourd'hui
(tableau~\ref{tab:poids_geles}). En dehors de la monnaie, l'erreur grandit quand l'échéance raccourcit : de
$-0{,}9\,\%$ à dix ans à $-3{,}0\,\%$ à un an sur le call à 110~\%, de $+1{,}4\,\%$ à $+4{,}3\,\%$ sur le put à
90~\%. La loi du prix de l'OAT, somme de lognormales, n'est pas lognormale : l'approximation sous-évalue sa queue
droite et surévalue sa queue gauche, et l'erreur relative pèse d'autant plus que l'option est loin de la monnaie
en nombre d'écarts-types, donc que l'échéance est courte (figure~\ref{fig:validation}).
\emph{L'hypothèse~H5 est confirmée à la monnaie et rejetée en dehors} ; le pricer utilise le prix exact.

\begin{figure}[htbp]
\centering
\includegraphics[width=\textwidth]{validation}
\caption{À gauche, call à cinq ans selon le strike : prix exact, poids gelés et Monte-Carlo ; à droite, erreur
relative des poids gelés selon le strike, à 1, 5 et 10 ans.}
\label{fig:validation}
\end{figure}

\subsection{Effet du facteur spread}
\label{sec:h3}

On compare le modèle à deux facteurs calibré au facteur taux seul ($\sigma_y = 0$, même forward repo, même
courbe), c'est-à-dire à ce que donnerait un pricer Hull-White calibré sur les seules swaptions €STR.

\begin{table}[htbp]
\centering\small
\caption{Effet du facteur spread sur le call à la monnaie forward (prix pour 100 de nominal)}
\label{tab:h3}
\begin{tabular}{@{}lccccccccc@{}}
\toprule
 & & \multicolumn{2}{c}{$\bar\Sigma_B$} & \multicolumn{3}{c}{Parts de $\bar\Sigma_B^2$} & \multicolumn{3}{c}{Call à la monnaie} \\
\cmidrule(lr){3-4}\cmidrule(lr){5-7}\cmidrule(lr){8-10}
$T_0$ & $F^{clean}$ & 2F & Taux seul & Taux & Spread & Croisée & 2F & Taux seul & Écart \\
\midrule
1 an   & 90,363 & 0,1007 & 0,0859 & 0,727 & 0,179 & 0,094 & 3,5781 & 3,0512 & +17,3~\% \\
2 ans  & 89,323 & 0,1373 & 0,1170 & 0,726 & 0,180 & 0,094 & 4,6770 & 3,9878 & +17,3~\% \\
5 ans  & 86,709 & 0,1909 & 0,1625 & 0,725 & 0,181 & 0,094 & 5,7627 & 4,9109 & +17,3~\% \\
7 ans  & 85,618 & 0,2030 & 0,1727 & 0,724 & 0,181 & 0,094 & 5,6794 & 4,8378 & +17,4~\% \\
10 ans & 85,266 & 0,1971 & 0,1677 & 0,724 & 0,182 & 0,094 & 4,9598 & 4,2223 & +17,5~\% \\
\bottomrule
\end{tabular}
\end{table}

Avec les paramètres estimés, le spread apporte 18~\% de la variance du prix forward de l'OAT, et le terme croisé
9~\% de plus, car $\rho > 0$ : le spread s'élargit quand les taux montent (tableau~\ref{tab:h3}). La volatilité
$\bar\Sigma_B$ dépasse de 17 à 18~\% celle du facteur taux seul, et le call à la monnaie est plus cher de 17~\% à
toutes les échéances (figure~\ref{fig:resultats}). \emph{L'hypothèse~H3 est confirmée.}

Le call à la monnaie croît jusqu'à cinq ans, puis décroît. Au-delà de cinq ans, l'obligation restante raccourcit :
sa volatilité de prix plafonne vers sept ans puis baisse, et le forward comme l'actualisation diminuent. Le forward
clean passe de 90,4 à un an à 85,3 à dix ans : porter une OAT qui rapporte plus que son financement fait baisser
son prix forward, et l'effet s'accumule avec l'échéance. Le recentrage $\delta$ vaut $-2{,}6$, $-3{,}2$, $-0{,}9$,
$8{,}3$ et $33{,}0$~pb aux cinq échéances : la courbe repo retenue est du même ordre que le spread de l'OAT, un peu
au-dessus à un an et en dessous à dix ans, et l'approximation que porte $\delta$ (section~\ref{sec:evaluation})
reste limitée.

\begin{figure}[htbp]
\centering
\includegraphics[width=\textwidth]{resultats}
\caption{À gauche, call à la monnaie forward selon l'échéance, modèle à deux facteurs et facteur taux seul ; à
droite, décomposition de $\bar\Sigma_B^2$ en parts taux, spread et croisée.}
\label{fig:resultats}
\end{figure}

\subsection{Un facteur suffit-il ?}
\label{sec:h4}

La répartition de la variance ne dépend presque pas de l'échéance, parce que les deux vitesses de retour sont
faibles : $H_{a_x} \approx H_{a_y}$, les deux facteurs déplacent les prix de la même façon, et leur somme se
comporte comme un facteur unique de volatilité
$\sigma_{eq} = \sqrt{\sigma_x^2 + \sigma_y^2 + 2\rho\sigma_x\sigma_y} = 0{,}849\,\%$, contre $\sigma_x = 0{,}737\,\%$.
On le vérifie avec un Hull-White à un facteur sur la courbe OAT, de vitesse $a_x$ et de volatilité $\sigma_{eq}$.

\begin{table}[htbp]
\centering\small
\caption{Hull-White à un facteur équivalent contre modèle à deux facteurs}
\label{tab:h4}
\begin{tabular}{@{}lccccc@{}}
\toprule
 & \multicolumn{3}{c}{Volatilité $\sigma_{eq} = 0{,}849\,\%$} & \multicolumn{2}{c}{Volatilité ajustée à la monnaie} \\
\cmidrule(lr){2-4}\cmidrule(lr){5-6}
$T_0$ & Call 2F & Call 1F & Écart à la monnaie & Volatilité (\%) & Écart max. 90-110~\% \\
\midrule
1 an   & 3,5781 & 3,5165 & $-1{,}72$~\% & 0,8641 & 0,053~\% \\
2 ans  & 4,6770 & 4,5951 & $-1{,}75$~\% & 0,8644 & 0,037~\% \\
5 ans  & 5,7627 & 5,6572 & $-1{,}83$~\% & 0,8651 & 0,020~\% \\
7 ans  & 5,6794 & 5,5729 & $-1{,}87$~\% & 0,8655 & 0,014~\% \\
10 ans & 4,9598 & 4,8643 & $-1{,}92$~\% & 0,8659 & 0,008~\% \\
\bottomrule
\end{tabular}
\end{table}

Le Hull-White à un facteur reproduit le call à la monnaie du modèle à deux facteurs à 1,9~\% près à toutes les
échéances, un peu en dessous, la sensibilité $H_{a_y}$ au spread dépassant légèrement $H_{a_x}$
(tableau~\ref{tab:h4}). En dehors de la monnaie, l'écart grandit jusqu'à $-5{,}1\,\%$ sur le put à 90~\% à un an,
mais il ne tient qu'au niveau de $\sigma_{eq}$, un peu trop bas. Avec une volatilité ajustée sur le call à la
monnaie, de 0,864 à 0,866~\%, le modèle à un facteur reproduit les prix du modèle à deux facteurs de 90 à 110~\% du
forward à 0,05~\% près. Le second facteur ne change donc pas la forme de la loi du prix de l'OAT à l'échéance,
seulement sa dispersion. Pour une option sur une seule OAT, il agit comme un supplément de volatilité ; sa
structure propre ne compterait que pour des produits sensibles à l'écart entre les deux courbes, comme les options
sur spread ou sur plusieurs titres. C'est la frontière que tracent \citet[section~4.1]{BrigoMercurio2006} entre
modèles à un et à deux facteurs, et \citet[section~6.4.6]{Chaix2021} de même : un facteur suffit quand les taux qui
font le payoff sont presque parfaitement corrélés. Avec des vitesses quasi nulles, chaque facteur est proche d'un
modèle de Ho et Lee, limite du Vasicek quand $a \to 0$ \citep[CM3, p.~14]{Hillairet2026}, et leur somme aussi.
\emph{L'hypothèse~H4 est rejetée avec les paramètres estimés.}

\subsection{Cohérence de la calibration mixte}

Le facteur taux est calibré sur des volatilités implicites, le facteur spread sur des volatilités historiques. Le
modèle calibré donne au swap €STR à vingt ans une volatilité de 66,3~pb par an et au rendement OAT à vingt ans
78,1~pb ; sur 2021-2026, les volatilités réalisées sont de 69,1 et 81,0~pb. Le rapport OAT~/~€STR, qui mesure
l'effet du spread, vaut 1,177 dans le modèle et 1,173 dans les données. Le contrôle n'est pas indépendant de la
calibration, puisque la volatilité du spread et $\rho$ viennent des mêmes séries, mais il montre que la volatilité
implicite des swaptions et les volatilités historiques ne se contredisent pas sur la période.

Comme un facteur d'actualisation exponentiel-affine laisse la variance inchangée, ce contrôle teste, pour les taux,
l'absence de prime de risque de variance que l'on doit supposer pour le spread. Le filtre de Kalman, estimé sur les
seules séries, permet de mener ce test sur toute la diagonale : avec sa volatilité historique des taux (0,775~\%),
il surévalue les volatilités des swaptions de 2,8~pb en moyenne (de 1,7 à 5,2~pb). La volatilité implicite à la
même vitesse de retour est de 0,744~\%, et le rapport $\sigma^{\Q}/\sigma^{\P}$ vaut 0,96. L'écart est
significatif, à 2,5 écarts-types, mais faible : s'il valait aussi pour le spread, il ne réduirait le call à la
monnaie que de 0,9~\%.

\subsection{Sensibilités}

Le tableau~\ref{tab:sensibilites} rassemble les variations du call à la monnaie dans chaque scénario, par rapport
au cas central.

\begin{table}[htbp]
\centering\small
\caption{Variation du call à la monnaie par rapport au cas central, en \%}
\label{tab:sensibilites}
\begin{tabular}{@{}llrrrrr@{}}
\toprule
Paramètre & Scénario & 1 an & 2 ans & 5 ans & 7 ans & 10 ans \\
\midrule
\multirow{2}{*}{Volatilité du spread} & $\sigma_y$ divisée par 2 & $-9{,}5$ & $-9{,}5$ & $-9{,}5$ & $-9{,}6$ & $-9{,}6$ \\
 & $\sigma_y$ multipliée par 2 & $+27{,}7$ & $+27{,}7$ & $+27{,}7$ & $+27{,}8$ & $+27{,}9$ \\
\multirow{2}{*}{Corrélation} & $\rho = 0{,}064$ & $-2{,}4$ & $-2{,}4$ & $-2{,}4$ & $-2{,}4$ & $-2{,}4$ \\
 & $\rho = 0{,}434$ & $+10{,}4$ & $+10{,}4$ & $+10{,}4$ & $+10{,}4$ & $+10{,}4$ \\
\multirow{2}{*}{Vitesse du spread} & $a_y = 0{,}215$ & $+4{,}3$ & $+2{,}6$ & $-0{,}7$ & $-1{,}7$ & $-1{,}9$ \\
 & $a_y = 1$ & $-0{,}3$ & $-5{,}2$ & $-9{,}6$ & $-10{,}3$ & $-10{,}2$ \\
Estimation du spread & filtre de Kalman & $+0{,}5$ & $+0{,}5$ & $+0{,}5$ & $+0{,}5$ & $+0{,}5$ \\
\multirow{2}{*}{Vitesse des taux} & $a_x = 0{,}05$ & $+2{,}9$ & $+2{,}3$ & $+0{,}8$ & $+0{,}1$ & $-0{,}4$ \\
 & $a_x = 0{,}2$ & $+14{,}9$ & $+9{,}8$ & $-0{,}1$ & $-3{,}0$ & $-3{,}2$ \\
\multirow{3}{*}{Régimes} & toujours calme & $-5{,}9$ & $-5{,}9$ & $-5{,}9$ & $-6{,}0$ & $-6{,}0$ \\
 & toujours agité & $+14{,}6$ & $+14{,}6$ & $+14{,}6$ & $+14{,}6$ & $+14{,}7$ \\
 & régimes simulés (hebdomadaires) & $-0{,}6$ & $-0{,}3$ & $-0{,}1$ & $-0{,}1$ & $-0{,}1$ \\
\multirow{3}{*}{Spread repo (fictif)} & courbe $-25$~pb & $-3{,}3$ & $-5{,}1$ & $-9{,}8$ & $-13{,}5$ & $-20{,}7$ \\
 & courbe $+25$~pb & $+3{,}4$ & $+5{,}2$ & $+10{,}5$ & $+14{,}8$ & $+23{,}9$ \\
 & financement au taux de l'OAT & $-0{,}6$ & $-0{,}6$ & $+3{,}1$ & $+10{,}6$ & $+27{,}1$ \\
\bottomrule
\end{tabular}
\end{table}

La corrélation est le paramètre le moins bien connu : sa fourchette historique, de 0,06 à 0,43, est bien plus
large que l'incertitude statistique de l'estimation, et le call y varie de $-2\,\%$ à $+10\,\%$. La vitesse de
retour du spread, à volatilité du spread à vingt ans inchangée puisque c'est elle que l'on observe, change le prix
de moins de 5~\% jusqu'à 0,215, la vitesse historique estimée par le filtre de Kalman. Seul un retour rapide, que
la structure par terme des volatilités exclut, le réduirait d'environ 10~\% aux longues échéances. Passer des
paramètres de la GMM à ceux du filtre ne change le call que de 0,5~\%.

La vitesse de retour des taux est imposée par la diagonale. Si on la fixe à 0,05 et qu'on recalibre $\sigma_x$, le
plus grand résidu de la diagonale passe de 2,3 à 3,8~pb, et le call change de 3~\% au plus. À $a_x = 0{,}2$, le
résidu atteint 11,1~pb et le call à un an monte de 15~\% : c'est l'erreur de calibration qui déplace le prix, pas
la vitesse de retour elle-même, qui ne pèse que sur les instruments différents de ceux de la calibration
\citep[section~6.3.8]{Chaix2021}.

Pour mesurer l'effet des régimes, on répartit entre eux la variance et la covariance moyennes de la GMM :
$\sigma_y$ vaut 0,24~\% en régime calme et 0,52~\% en régime agité, avec une corrélation commune de 0,14. Un
spread qui resterait toute la vie de l'option dans le régime calme baisserait le call de 6~\%, dans le régime
agité l'augmenterait de 15~\%. Mais les régimes ne durent que 43 et 16 jours ouvrés en moyenne : sur un an ou plus,
le spread passe dans chacun une part de son temps proche de sa probabilité stationnaire, et, les vitesses de
retour étant quasi nulles, seule compte la variance cumulée. En moyenne sur 20\,000 trajectoires de régimes
simulées depuis la probabilité du régime agité au 8 septembre 2026, le call s'écarte de celui du modèle retenu de
0,2~\% à un an ; avec les régimes plus longs estimés sur variations hebdomadaires, de 0,6~\% au plus. Sous la
mesure de pricing, chaque régime garde sa variance \citep[cours~2, prop.~III.2, slide~69]{MPR2025} ; seules les
probabilités de transition pourraient différer, si le risque de régime était rémunéré, ce que l'on exclut.

Reste le financement. À un an, le niveau du repo pèse peu : translater la courbe de $-25$ à $+25$~pb déplace le
call de $-3\,\%$ à $+3\,\%$. À dix ans, la même translation fait passer ce call de 3,93 à 6,15, autour de 4,96
($-21\,\%$ à $+24\,\%$), et un financement au taux de l'OAT elle-même le porterait à 6,31
(figure~\ref{fig:sensibilites}). Aux longues échéances, le niveau du financement, fictif ici, pèse donc autant que
le facteur spread lui-même.

\begin{figure}[htbp]
\centering
\includegraphics[width=\textwidth]{sensibilites}
\caption{À gauche, forward clean de l'OAT 2046 selon la courbe repo ; à droite, call à la monnaie selon la
volatilité du spread.}
\label{fig:sensibilites}
\end{figure}
